\subsubsection{Barra articulada a argolla circular}

Una barra (con espesor mayor que cero) de longitud $\ell$ está articulada en uno de sus extremos a una argolla, como se muestra en la Figura \ref{fig:barra_anillo}. La argolla se desliza a lo largo de una pista fija en forma de anillo circular de radio $R$. Sabiendo que el ángulo $\theta$ describe la posición de la argolla con respecto al centro del anillo y que el ángulo $\phi$ describe la orientación de la barra con respecto a la línea radial instantánea desde el centro del anillo hasta la argolla, determine las siguientes cantidades vistas por un observador fijo a la pista:

\begin{enumerate}
    \item Velocidad angular de la barra.
    \item Velocidad y aceleración del extremo libre de la barra (punto $P$).
\end{enumerate}

\begin{figure}[!ht]
    \centering
    \begin{tikzpicture}[>=Latex, scale=1.8]
        % --- Parámetros ---
        \def\R{1.5}      % Radio del anillo
        \def\L{1.8}      % Longitud de la barra l
        \def\ThetaAng{40} % Ángulo theta
        \def\PhiAng{35}   % Ángulo phi

        % --- 1. Soportes del Anillo ---
        \foreach \s in {1, -1}{
            \begin{scope}[yshift=\s*\R cm]
                \draw[pattern=north east lines, pattern color=gray!60] (-0.2, 0) rectangle (0.2, 0.15*\s);
                \draw[thick] (-0.3, 0.15*\s) -- (0.3, 0.15*\s);
                \draw[thick] (0,0) -- (0, 0.15*\s);
            \end{scope}
        }

        % --- 2. Anillo Circular ---
        \draw[thick, double distance=2pt] (0,0) circle (\R);
        \coordinate (O) at (0,0);
        \filldraw (O) circle (1.2pt) node[left] {$O$};

        % --- 3. Argolla y Barra ---
        % Posición de la argolla Q
        \coordinate (Q) at ({-90+\ThetaAng}:\R);
        % Posición del extremo P
        \coordinate (P) at ($(Q) + ({-90+\ThetaAng+\PhiAng}:\L)$);

        % Dibujo de la barra con grosor
        \draw[line width=5pt, line cap=round, color=black] (Q) -- (P);
        \draw[line width=4pt, line cap=round, color=gray!30] (Q) -- (P);
        \node[below] at ($(Q)!0.5!(P)$) {$\ell$};

        % --- 4. Líneas de Referencia ---
        \draw[dashed, gray] (O) -- (0, -\R-0.4); % Vertical
        \draw[dashed] (O) -- (Q);                % Radio OQ
        \draw[dotted, thick] (Q) -- ($(O)!1.7!(Q)$); % Extensión radial
        
        % Flecha del radio R
        \draw[<->] (O) -- (30:\R) node[midway, above] {$R$};

        % --- 5. Ángulos ---
        % Ángulo theta
        \draw[->] (0, -0.7) arc (-90:{-90+\ThetaAng}:0.7);
        \node at (-75:0.55) {$\theta$};

        % Ángulo phi
        \begin{scope}[shift={(Q)}]
            \draw[->] ({-90+\ThetaAng}:0.5) arc ({-90+\ThetaAng}:{-90+\ThetaAng+\PhiAng}:0.5);
            \node at ({-90+\ThetaAng+0.5*\PhiAng}:0.7) {$\phi$};
        \end{scope}

        % --- 6. Etiquetas de puntos ---
        \filldraw[fill=white, thick] (Q) circle (2.5pt); % Argolla
        \filldraw (Q) circle (1pt) node[below] {$Q$};
        \filldraw (P) circle (1.2pt) node[right] {$P$};

    \end{tikzpicture}
    \caption{Barra de longitud $\ell$ articulada a una argolla $Q$ que se desplaza sobre una guía circular.}
    \label{fig:barra_anillo}
\end{figure}